\clearpage
\section{Universal lane transport at finite profile}
\label{sec:block15-universal-lane-transport}

The preceding parts of the manuscript establish packet coordinates in an
additive idempotent-complete setting and verify them separately in several
native strata.  This section isolates the common finite theorem.  The word
``lane'' will mean a packet together with its address, its chosen splitting,
and its position in a declared refinement system.  Thus a lane is more than a
subsequence or a suggestive decomposition.  It carries the exact data needed
for analysis, synthesis, transfer, refinement, and reconstruction.

The theorem is universal only in a profiled sense.  Its algebraic core applies
to every additive realization that preserves the finite sums, compositions,
split idempotents, and refinement identities used in the proof.  Kernel,
quotient, tensor, topological, metric, and nonlinear conclusions require the
additional hypotheses stated where they occur.

\subsection{Finite profiled lane systems}
\label{subsec:block15-lane-systems}

\begin{definition}[Finite resolution shape]
\label{def:block15-resolution-diagram}
A finite resolution shape is a finite partially ordered set
\(\mathsf R\).  We write \(s\succeq r\) when \(s\) is finer than \(r\).
The shape contains no address sets: object-specific addresses and their
refinement maps are declared in the lane system below.
\end{definition}

\begin{definition}[Finite profiled lane system]
\label{def:block15-finite-profiled-lane-system}
Let \(\mathcal C\) be an additive idempotent-complete category, let
\(X\in\mathcal C\), and let \(\mathsf R\) be a finite resolution shape.  A
finite profiled lane system \(\mathcal L_X\) consists, for every
\(r\in\mathsf R\), of a finite address set \(I_{X,r}\) and a family
\[
 \mathcal P_X^r=(p_{X,i}^r)_{i\in I_{X,r}}
 \subseteq\End_{\mathcal C}(X)
\]
with
\[
 (p_{X,i}^r)^2=p_{X,i}^r,
 \qquad
 p_{X,i}^rp_{X,j}^r=0\quad(i\ne j),
 \qquad
 \sum_{i\in I_{X,r}}p_{X,i}^r=1_X,
\]
together with chosen splittings
\[
 X_i^r\xrightarrow{\iota_{X,i}^r}X
 \xrightarrow{\pi_{X,i}^r}X_i^r,
 \qquad
 \iota_{X,i}^r\pi_{X,i}^r=p_{X,i}^r,
 \qquad
 \pi_{X,i}^r\iota_{X,i}^r=1_{X_i^r}.
\]
For \(s\succeq r\), the system also contains a surjection
\[
 \rho^X_{sr}:I_{X,s}\twoheadrightarrow I_{X,r}
\]
such that
\[
 p_{X,i}^r
 =
 \sum_{\rho^X_{sr}(j)=i}p_{X,j}^s,
 \qquad
 \rho^X_{tr}=\rho^X_{sr}\rho^X_{ts}.
\]
Here \(\rho^X_{rr}=1_{I_{X,r}}\).  The address sets and refinement maps are
declared data; they are not inferred from their cardinalities.
Its verified profile contains finite biproducts, addition of morphisms,
composition, the displayed splittings, and these projector-sum refinement
identities.  Any further constructor is included only after its preservation
has been verified.
\end{definition}

At resolution \(r\), write
\[
 \Analysis_X^r
 :=(\pi_{X,i}^r)_i:
 X\longrightarrow\bigoplus_iX_i^r,
 \qquad
 \Synthesis_X^r
 :=(\iota_{X,i}^r)_i:
 \bigoplus_iX_i^r\longrightarrow X.
\]
For profiled lane systems on \(X\) and \(Y\), define the coordinate block of
\(f:X\to Y\) at resolutions \(r,u\) by
\[
 [f]^{u,r}_{a i}
 :=
 \pi_{Y,a}^u f\iota_{X,i}^r:
 X_i^r\longrightarrow Y_a^u.
\]
Its ambient corner is
\[
 p_{Y,a}^u f p_{X,i}^r
 =
 \iota_{Y,a}^u[f]^{u,r}_{a i}\pi_{X,i}^r.
\]

\subsection{The universal finite transport theorem}
\label{subsec:block15-main-theorem}

\begin{theorem}[Universal Lane Transport Theorem, finite profiled form]
\label{thm:block15-universal-lane-transport}
Let \(\mathcal C\) be an additive idempotent-complete category.  Give each of
the objects \(X_0,\ldots,X_m\) a finite profiled lane system.  Then all of the
following statements hold.

\begin{enumerate}[label=\textup{(L\arabic*)}]
\item \textbf{Exact analysis and synthesis.}  At every declared resolution,
\[
 \Synthesis_X^r\Analysis_X^r=1_X,
 \qquad
 \Analysis_X^r\Synthesis_X^r
 =1_{\oplus_iX_i^r}.
\]
Hence \(X\cong\bigoplus_iX_i^r\).

\item \textbf{Unique lane matrix.}  Every morphism \(f:X\to Y\) is recovered
uniquely from its lane blocks:
\[
 f
 =
 \sum_{a,i}p_{Y,a}^u f p_{X,i}^r
 =
 \sum_{a,i}\iota_{Y,a}^u
 [f]^{u,r}_{a i}\pi_{X,i}^r.
\]
Equivalently, the assignment
\[
 f\longmapsto
 \bigl([f]^{u,r}_{a i}\bigr)_{a,i}
\]
is a bijection from \(\Hom_{\mathcal C}(X,Y)\) to the group of rectangular
matrices with entries in
\(\Hom_{\mathcal C}(X_i^r,Y_a^u)\).

\item \textbf{Composition and path law.}  If \(f:X\to Y\) and \(g:Y\to Z\),
with the same resolution \(u\) used on the intermediate object, then
\[
 [gf]^{v,r}_{c i}
 =
 \sum_{a\in I_{Y,u}}[g]^{v,u}_{c a}[f]^{u,r}_{a i}.
\]
More generally, a block of a composable string is the finite sum of the
products along all compatible intermediate address paths.

\item \textbf{Projector-sum refinement.}  If \(s\succeq r\) on \(X\) and
\(v\succeq u\) on \(Y\), then
\[
 p_{Y,a}^u f p_{X,i}^r
 =
 \sum_{\rho^Y_{vu}(b)=a}
 \sum_{\rho^X_{sr}(j)=i}
 p_{Y,b}^v f p_{X,j}^s.
\]
Coarsening a fine composite gives the same ambient morphism as composing at
the coarse level.

\item \textbf{Selectors and named residuals.}  For
\(E\subseteq I_{Y,u}\times I_{X,r}\), put
\[
 C_E^{u,r}(f)
 :=\sum_{(a,i)\in E}p_{Y,a}^ufp_{X,i}^r,
 \qquad
 R_E^{u,r}(f):=f-C_E^{u,r}(f).
\]
As additive endomorphisms of \(\Hom_{\mathcal C}(X,Y)\),
\[
 C_E^2=C_E,
 \qquad
 C_EC_F=C_{E\cap F},
 \qquad
 R_E=1-C_E,
 \qquad
 f=C_E(f)+R_E(f).
\]
Thus deletion of a lane block is never silent.  Its complement is an exact,
named residual.

\item \textbf{Support transport.}  If \(f=C_E(f)\) and \(g=C_F(g)\), then
\[
 gf=C_{F\circ E}(gf),
\]
where \(F\circ E\) is relational composition through the intermediate
addresses.  For arbitrary \(f,g\), writing
\(f_E=C_E(f)\), \(r_f=R_E(f)\), \(g_F=C_F(g)\), and
\(r_g=R_F(g)\), one has the exact identity
\[
 gf-g_Ff_E
 =r_gf_E+g_Fr_f+r_gr_f.
\]

\item \textbf{Additive realization.}  Let
\(F:\mathcal C\to\mathcal D\) be an additive functor to another additive
idempotent-complete category.  If the verified profile of \(F\) contains the
chosen splittings and refinement identities, then
\[
 F([f]^{u,r}_{a i})
 =
 [F(f)]^{u,r}_{a i}
\]
after the declared splitting comparisons.  The functor carries analysis,
synthesis, matrix multiplication, projector-sum refinement, selectors, and
named residuals to their corresponding constructions in \(\mathcal D\).

\item \textbf{Exact fold loss.}  Let \(E\) be a finite set of retained lane
positions and set
\[
 B_e:=\Hom_{\mathcal C}(X_i^r,Y_a^u)
 \quad\text{for }e=(a,i)\in E.
\]
Let \(T\) be a finite set, let \(\kappa:E\to T\), let \(C_t\) be abelian
groups, and let \(\tau_e:B_e\to C_{\kappa(e)}\) be homomorphisms.  With
\(C=\bigoplus_{t\in T}C_t\) and coordinate inclusions
\(j_t:C_t\to C\), the fold
\[
 \Phi_{\kappa,\tau}:
 \bigoplus_{e\in E}B_e\longrightarrow C,
 \qquad
 (b_e)_e\longmapsto
 \sum_{e\in E}j_{\kappa(e)}\tau_e(b_e)
\]
has exact information-loss group \(\ker\Phi_{\kappa,\tau}\).  In particular,
two retained lane matrices have the same folded record exactly when their
difference belongs to this kernel.

Suppose, more strongly, that every \(\tau_e\) is injective and that there are
globally typed selectors \(Q_e\in\End(C)\) satisfying, for every
\(e,e'\in E\) and every \(b_{e'}\in B_{e'}\),
\[
 Q_ej_{\kappa(e')}\tau_{e'}(b_{e'})
 =
 \delta_{e,e'}j_{\kappa(e')}\tau_{e'}(b_{e'}).
\]
Then \(\Phi_{\kappa,\tau}\) is injective.  If there are also homomorphisms
\[
 r_e:C\longrightarrow B_e,
 \qquad
 r_ej_{\kappa(e)}\tau_e=1_{B_e},
\]
then
\[
 L:=\bigl(r_eQ_e\bigr)_{e\in E}:C\longrightarrow\bigoplus_{e\in E}B_e
\]
is a left inverse of \(\Phi_{\kappa,\tau}\), so the fold is split lossless.
\end{enumerate}
\end{theorem}

\begin{proof}
For (L1), the first composite is
\[
 \sum_i\iota_{X,i}^r\pi_{X,i}^r
 =\sum_i p_{X,i}^r=1_X.
\]
The \((i,j)\)-entry of the other composite is
\(\pi_{X,i}^r\iota_{X,j}^r\).  It is the identity when \(i=j\), and it is
zero otherwise because the corresponding idempotents annihilate one another.

For (L2), insert the two reconstruction identities on the two sides of
\(f\).  Applying \(\pi_{Y,a}^u\) on the left and
\(\iota_{X,i}^r\) on the right proves uniqueness of every entry.  For (L3),
insert
\(1_Y=\sum_a\iota_{Y,a}^u\pi_{Y,a}^u\) between \(g\) and \(f\).
Iteration gives the path formula.

For (L4), substitute the two projector-sum refinement identities into the
coarse ambient corner and distribute the finite sums.  The assertion about
composition follows from the same substitution at the intermediate object.

For (L5), each selector either fixes an ambient block or kills it.  Applying
two selectors therefore retains precisely the intersection.  The residual
identities follow by subtraction.  For (L6), expand each block of \(gf\)
through the intermediate resolution.  A surviving summand supplies an
intermediate address in the relational product.  The residual formula is the
finite expansion of
\((g_F+r_g)(f_E+r_f)-g_Ff_E\).

For (L7), additivity preserves zero maps and finite sums, functoriality
preserves identities and composition, and the splitting comparisons identify
the image of every source lane object with the declared target lane object.
The profile hypothesis is needed precisely for those identifications and for
the refinement equations.

For (L8), equality of two fold values is equivalent to the difference of the
two input families lying in the kernel.  Under the selector hypothesis, if
\(\Phi_{\kappa,\tau}((b_e)_e)=0\), applying \(Q_e\) gives
\(j_{\kappa(e)}\tau_e(b_e)=0\).  Both \(j_{\kappa(e)}\) and \(\tau_e\) are
injective, hence \(b_e=0\) for every \(e\).  The chosen retractions compose
with the \(Q_e\) to give the displayed left inverse: its \(e\)-coordinate
on the fold is \(r_eQ_ej_{\kappa(e)}\tau_e=1_{B_e}\), while strict
separation kills every other coordinate.
\end{proof}

\begin{remark}[Why the fold selectors live in one group]
\label{rem:block15-global-fold-typing}
The endomorphisms \(Q_e\) in (L8) all act on the single group
\(C=\bigoplus_tC_t\).  Writing separate selectors on unrelated codomains
would make the Kronecker relation ill-typed and would not prove injectivity of
the combined fold.
\end{remark}

\begin{proposition}[Exact visible-support loss ledger]
\label{prop:block15-visible-support-ledger}
Use the finite fold data of (L8) and fix an input family
\(b=(b_e)_{e\in E}\).  Define
\[
 S_0:=\{e\in E:b_e\ne0\},
 \qquad
 A_0:=\{e\in S_0:\tau_e(b_e)=0\},
 \qquad
 S_1:=S_0\setminus A_0,
\]
and let
\[
 \Gamma:=\kappa(S_1)\subseteq T.
\]
For \(t\in T\), write the folded coordinate as
\[
 \overline b_t
 :=
 \sum_{\substack{e\in S_1\\\kappa(e)=t}}\tau_e(b_e)
 \in C_t,
\]
and put
\[
 \operatorname{supp}(\overline b)
 :=\{t\in T:\overline b_t\ne0\},
 \qquad
 K_0:=\{t\in\Gamma:\overline b_t=0\}.
\]
Then
\[
 \boxed{
 |S_0|-|\operatorname{supp}(\overline b)|
 =
 |A_0|
 +\bigl(|S_1|-|\Gamma|\bigr)
 +|K_0|.}
\]
The three terms record, respectively, annihilation by the local maps
\(\tau_e\), collision of surviving lane addresses under \(\kappa\), and
cancellation among folded values at a common address.  If every \(\tau_e\) is
injective, then \(A_0=\varnothing\).  If globally typed selectors satisfy the
strict separation law of (L8), then \(K_0=\varnothing\).
\end{proposition}

\begin{proof}
The first partition gives \(|S_0|=|A_0|+|S_1|\).  The definition of
\(K_0\) gives the disjoint partition
\[
 \Gamma
 =\operatorname{supp}(\overline b)\sqcup K_0,
\]
so \(|\Gamma|=|\operatorname{supp}(\overline b)|+|K_0|\).  Subtracting the
second equality from the first proves the formula.

Injectivity of \(\tau_e\) makes \(\tau_e(b_e)\ne0\) for every \(e\in S_0\),
so there is no annihilation.  For strict selectors, suppose
\(t\in K_0\) and choose \(e\in S_1\) with \(\kappa(e)=t\).  Applying \(Q_e\)
to
\[
 j_t\overline b_t
 =\sum_{\kappa(e')=t}j_t\tau_{e'}(b_{e'})=0
\]
leaves \(j_t\tau_e(b_e)=0\), contrary to \(e\in S_1\).  Hence
\(K_0=\varnothing\).
\end{proof}

\subsection{Certified change of resolution and exact loss sequences}
\label{subsec:block15-coarsening-loss}

For resolutions \(s\succeq r\) on \(X\) and \(v\succeq u\) on \(Y\), let
\[
 \mathsf B^{v,s}(X,Y)
 :=
 \bigoplus_{b\in I_{Y,v}}
 \bigoplus_{j\in I_{X,s}}
 \Hom_{\mathcal C}(X_j^s,Y_b^v),
\]
and define \(\mathsf B^{u,r}(X,Y)\) similarly.  Analysis of a morphism gives
an isomorphism
\[
 \mathfrak a^{v,s}_{X,Y}:
 \Hom_{\mathcal C}(X,Y)\xrightarrow{\sim}
 \mathsf B^{v,s}(X,Y),
\]
whose inverse is lane synthesis.

\begin{proposition}[Certified coarsening is an exact coordinate change]
\label{prop:block15-certified-coarsening}
The change from fine to coarse block coordinates, canonical relative to the
chosen splittings, is
\[
 \Gamma^{(u,r)}_{(v,s)}
 :=
 \mathfrak a^{u,r}_{X,Y}
 \bigl(\mathfrak a^{v,s}_{X,Y}\bigr)^{-1}.
\]
For a fine matrix \(h=(h_{bj})\), its \((a,i)\)-entry is
\[
 \bigl(\Gamma^{(u,r)}_{(v,s)}h\bigr)_{a i}
 =
 \sum_{b,j}
 \bigl(\pi_{Y,a}^u\iota_{Y,b}^v\bigr)
 h_{bj}
 \bigl(\pi_{X,j}^s\iota_{X,i}^r\bigr).
\]
Only descendants of \(a\) and \(i\) can contribute.  The inverse is the
corresponding fine analysis after coarse synthesis.  Consequently, changing a
complete lane resolution loses no information.  Loss occurs only after a
subsequent selector, quotient, summary, or fold is applied.
\end{proposition}

\begin{proof}
The displayed definition is a composite of two isomorphisms from
\cref{thm:block15-universal-lane-transport}.  Expanding coarse analysis after
fine synthesis gives the entry formula.  If \(b\) is not a descendant of
\(a\), then
\(p_{Y,a}^up_{Y,b}^v=0\), hence
\(\pi_{Y,a}^u\iota_{Y,b}^v=0\).  The source factor vanishes by the same
argument when \(j\) is not a descendant of \(i\).  Reversing analysis and
synthesis produces the inverse.
\end{proof}

\begin{corollary}[Functorial coarsening]
\label{cor:block15-functorial-coarsening}
For \(t\succeq s\succeq r\) and compatible target resolutions,
the certified coordinate changes compose strictly.  They also intertwine
matrix multiplication.  Thus a finite refinement diagram gives a genuine
diagram of equivalent matrix-coordinate categories, not merely a family of
numerically similar arrays.
\end{corollary}

\begin{theorem}[Abelian fold-loss exact sequence]
\label{thm:block15-abelian-fold-exact-sequence}
Let \(\mathcal A\) be an abelian category.  Suppose a finite family of lane
objects \((B_e)_{e\in E}\), target objects \((C_t)_{t\in T}\), and morphisms
\(\tau_e:B_e\to C_{\kappa(e)}\) is given.  Let
\(q_e:\bigoplus_{e'\in E}B_{e'}\to B_e\) be the coordinate projection and
define
\[
 \Phi:=\sum_{e\in E}j_{\kappa(e)}\tau_eq_e:
 \bigoplus_{e\in E}B_e\longrightarrow
 \bigoplus_{t\in T}C_t.
\]
Then there are canonical short exact sequences
\[
 0\longrightarrow\ker\Phi
 \longrightarrow\bigoplus_eB_e
 \longrightarrow\im\Phi
 \longrightarrow0
\]
and
\[
 0\longrightarrow\im\Phi
 \longrightarrow\bigoplus_tC_t
 \longrightarrow\coker\Phi
 \longrightarrow0.
\]
The first sequence is the exact fold-loss certificate.  Suppose every
\(\tau_e\) is monic and there are endomorphisms
\(Q_e\in\End(\bigoplus_tC_t)\) satisfying
\[
 Q_ej_{\kappa(e')}\tau_{e'}
 =\delta_{e,e'}j_{\kappa(e')}\tau_{e'}
 \quad\text{as morphisms }B_{e'}\longrightarrow\bigoplus_tC_t.
\]
Then \(\Phi\) is monic.  If, in addition, there are morphisms
\[
 r_e:\bigoplus_tC_t\longrightarrow B_e,
 \qquad r_ej_{\kappa(e)}\tau_e=1_{B_e},
\]
then \(L=(r_eQ_e)_e\) is a left inverse of \(\Phi\).  Thus the first exact
sequence splits, \(\ker\Phi=0\), and the image inclusion in the second exact
sequence also splits.
\end{theorem}

\begin{proof}
The two sequences are the kernel--image and image--cokernel factorizations in
an abelian category.  The selector argument in (L8), interpreted internally
in \(\mathcal A\), makes \(\Phi\) monic.  The displayed morphism \(L\) is a
left inverse by strict separation.  Hence the first sequence splits;
moreover, \(\Phi L\) is an idempotent with image \(\im\Phi\), so the image
inclusion in the second sequence splits as well.
\end{proof}

\begin{boundary}[Coordinate coarsening versus statistic coarsening]
\label{bnd:block15-coordinate-vs-statistic-coarsening}
The complete fine and coarse block matrices are isomorphic coordinate records.
Replacing each coarse block by a dimension, norm, valuation, degree, support
bit, or color is a different map.  Its kernel or fiber must be computed in the
native doctrine.  The lane theorem does not make such a statistic complete.
\end{boundary}

\subsection{The matrix category carried by a lane system}
\label{subsec:block15-matrix-category}

\begin{definition}[Lane matrix category at a fixed resolution]
\label{def:block15-lane-matrix-category}
Let \(\mathcal L_{\mathsf R}\) be the category whose objects are objects of
\(\mathcal C\) equipped with full finite profiled lane systems over
\(\mathsf R\), with all underlying \(\mathcal C\)-morphisms allowed.  For
\(r\in\mathsf R\), define
\(\operatorname{Mat}_r(\mathcal L_{\mathsf R})\) to have the tagged finite
families
\([X;(X_i^r)_{i\in I_{X,r}}]\) as objects and rectangular matrices
\[
 (h_{a i}),
 \qquad h_{a i}:X_i^r\to Y_a^r,
\]
as morphisms.  Composition is matrix contraction and identities are diagonal
identity matrices.

For \(s\succeq r\), define
\[
 \Gamma_{sr}:\operatorname{Mat}_s(\mathcal L_{\mathsf R})
 \longrightarrow\operatorname{Mat}_r(\mathcal L_{\mathsf R})
\]
by replacing each object tag with its \(r\)-tag and sending a fine matrix to
the certified source-and-target coordinate change of
\cref{prop:block15-certified-coarsening}.
\end{definition}

\begin{theorem}[Fixed-resolution matrix-category equivalence]
\label{thm:block15-matrix-category-equivalence}
The functor
\[
 \mathfrak M_r:\mathcal L_{\mathsf R}\longrightarrow
 \operatorname{Mat}_r(\mathcal L_{\mathsf R}),
 \qquad
 X\longmapsto[X;(X_i^r)_i],
 \qquad
 f\longmapsto([f]_{a i}^{r,r}),
\]
is an equivalence.  A quasi-inverse returns the lane-decorated object named by
the tag and sends a matrix to
\[
 \sum_{a,i}\iota_{Y,a}^r h_{a i}\pi_{X,i}^r.
\]
For every \(s\succeq r\), the coordinate changes are functors and satisfy
\[
 \Gamma_{sr}\mathfrak M_s=\mathfrak M_r.
\]
For \(t\succeq s\succeq r\), they also satisfy
\[
 \Gamma_{tr}=\Gamma_{sr}\Gamma_{ts}.
\]
Thus the equivalences are compatible with every certified change of
resolution.
\end{theorem}

\begin{proof}
The unique lane matrix and composition clauses of
\cref{thm:block15-universal-lane-transport} show that \(\mathfrak M_r\) is
fully faithful.  Every object of the displayed matrix category is, by
definition, the lane family of an object of \(\mathcal L_{\mathsf R}\), so it is
essentially surjective.  Analysis and synthesis furnish the unit and counit
isomorphisms.  Each \(\Gamma_{sr}\) is conjugation by exact analysis and
synthesis, so it preserves identities and matrix composition.  The two
displayed identities follow from cancellation of adjacent analysis--synthesis
pairs, equivalently from functorial coarsening in
\cref{cor:block15-functorial-coarsening}.
\end{proof}

\begin{remark}[No dimension claim]
\label{rem:block15-address-count-not-dimension}
The number \(|I_{X,r}|\) counts declared addresses.  A lane object can have
arbitrary dimension, rank, length, measure, or entropy, including zero where
the ambient doctrine permits it.  Matrix equivalence therefore does not
identify address count with any of those invariants.
\end{remark}

\subsection{Multilinear and specified-polynomial transport}
\label{subsec:block15-multilinear}

The additive theorem does not automatically govern arbitrary nonlinear maps.
There is, however, an exact extension in the category of modules when the
multilinear pieces are supplied as data.

\begin{theorem}[Module multilinear lane expansion]
\label{thm:block15-multilinear-lane-expansion}
Let \(\kk\) be a commutative ring.  Let
\[
 X^{(1)},\ldots,X^{(d)},Y
\]
be \(\kk\)-modules with finite split lane systems, and let
\[
 T:X^{(1)}\times\cdots\times X^{(d)}\longrightarrow Y
\]
be a specified \(\kk\)-multilinear map.  Define
\[
 T_{a;i_1,\ldots,i_d}
 :=
 \pi_{Y,a}T
 (\iota_{1,i_1}-,\ldots,\iota_{d,i_d}-).
\]
Then
\[
 T(x_1,\ldots,x_d)
 =
 \sum_{a,i_1,\ldots,i_d}
 \iota_{Y,a}
 T_{a;i_1,\ldots,i_d}
 (\pi_{1,i_1}x_1,\ldots,\pi_{d,i_d}x_d).
\]
Substitution of multilinear maps is governed by finite contraction over every
intermediate lane address.

Separately, let \(X,Y\) be \(\kk\)-modules with finite split lane systems and
suppose a polynomial map is presented with specified data
\[
 P(x)=T_0+T_1(x)+T_2(x,x)+\cdots+T_D(x,\ldots,x),
\]
where
\[
 T_0\in Y,
 \qquad
 T_n:X^n\longrightarrow Y\quad(1\le n\le D)
\]
and every \(T_n\) is a supplied multilinear map.  Then
\[
 T_0=\sum_a\iota_{Y,a}\pi_{Y,a}(T_0),
\]
and adding this degree-zero decomposition to the displayed multilinear
expansions gives an exact lane expansion of \(P\).
\end{theorem}

\begin{proof}
Insert \(x_\ell=\sum_{i_\ell}p_{\ell,i_\ell}x_\ell\) in every input, use
multilinearity to distribute the finite sums, and insert
\(1_Y=\sum_a p_{Y,a}\) in the output.  The coordinate formula follows from
the chosen splittings.  Repeating the calculation after substitution gives
the contraction law.  The degree-zero identity is output synthesis applied
to \(T_0\); the polynomial statement is then a finite sum of the specified
homogeneous pieces.
\end{proof}

\begin{boundary}[The nonlinear doctrine must be specified]
\label{bnd:block15-multilinear-boundary}
Over a general commutative ring, a polynomial function need not determine a
unique symmetric multilinear polarization.  Infinite analytic expansions also
require a declared topology and convergence theorem.  Therefore
\cref{thm:block15-multilinear-lane-expansion} is stated for supplied
multilinear maps and supplied finite polynomial data.  It does not assert a
canonical lane calculus for every nonlinear category.
\end{boundary}

\subsection{Factorial lanes and residue packets}
\label{subsec:block15-factorial-realization}

The factorial construction is an exact arithmetic realization of the abstract
lane mechanism.  It is useful because the addresses, refinement maps, Fourier
projectors, and hypergeometric coordinates can all be written explicitly.

For integers \(N\ge1\) and \(n\ge0\), define the step factorial
\[
 (a)_{n;N}:=\prod_{k=0}^{n-1}(a+Nk),
 \qquad (a)_{0;N}:=1.
\]
For a commutative ring \(A\) and a formal series
\(F(z)=\sum_{m\ge0}a_mz^m\in A[[z]]\), define its residue lane
\[
 \mathcal P_{N,j}F(z)
 :=
 \sum_{\substack{m\ge0\\m\equiv j\ (\mathrm{mod}\ N)}}a_mz^m,
 \qquad 0\le j<N.
\]

\begin{theorem}[Exact factorial--residue realization]
\label{thm:block15-factorial-residue-realization}
Let \(N,c,n\in\mathbb Z\) with \(N\ge1\), \(c\ge0\), and \(n\ge0\).
Then
\[
 \boxed{
 \frac{(Nn+c)!}{c!}
 =
 \prod_{r=c+1}^{c+N}(r)_{n;N}.}
\]
Each integer factor from \(c+1\) through \(c+Nn\) occurs in exactly one
lane.  For every commutative ring \(A\), on the \(A\)-module \(A[[z]]\) the
operators
\((\mathcal P_{N,j})_{0\le j<N}\) satisfy
\[
 \mathcal P_{N,j}^2=\mathcal P_{N,j},
 \qquad
 \mathcal P_{N,j}\mathcal P_{N,k}=0\ (j\ne k),
 \qquad
 \sum_{j=0}^{N-1}\mathcal P_{N,j}=1.
\]
If \(M\in\mathbb Z\) and \(M\ge1\), their exact refinement law is
\[
 \mathcal P_{N,j}
 =
 \sum_{\ell=0}^{M-1}\mathcal P_{NM,j+\ell N}.
\]

Suppose further that \(K\) is a field, that \(N\) is invertible in \(K\),
and that \(K\) contains a primitive \(N\)-th root of unity \(\omega\).  Then
for every \(F\in K[[z]]\),
\[
 \mathcal P_{N,j}F(z)
 =
 \frac1N\sum_{\ell=0}^{N-1}
 \omega^{-j\ell}F(\omega^\ell z).
\]
If, in addition, \(K\) has characteristic zero, the exponential and
hypergeometric series below are defined in \(K[[z]]\).  Put
\[
 E_{N,j}(z)
 :=\mathcal P_{N,j}e^z
 =\sum_{n\ge0}\frac{z^{Nn+j}}{(Nn+j)!}.
\]
If
\[
 \mathbf b_{N,j}
 :=
 \left\{
 \frac{j+r}{N}:1\le r\le N,
 \ r\ne N-j
 \right\},
\]
then
\[
 E_{N,j}(z)
 =
 \frac{z^j}{j!}
 {}_0F_{N-1}
 \left(
 \begin{matrix}-\\ \mathbf b_{N,j}\end{matrix};
 \frac{z^N}{N^N}
 \right).
\]
Thus coefficient-residue, root-of-unity, and hypergeometric descriptions are
exact additive coordinate realizations of one residue-address decomposition.
The step-factorial identity is an exact multiplicative companion with the
same address arithmetic; it is not asserted to be an additive idempotent
decomposition of the factorial.
\end{theorem}

\begin{proof}
Every integer \(m\) with \(c<m\le c+Nn\) has a unique representation
\(m=r+Nk\) with \(c+1\le r\le c+N\) and \(0\le k<n\).  Multiplying these
disjoint factors proves the factorial identity.

The projector laws follow coefficient by coefficient.  The residue classes
modulo \(NM\) lying over \(j\pmod N\) are precisely
\(j,j+N,\ldots,j+(M-1)N\), which proves refinement.  Character
orthogonality gives the Fourier formula.  Finally,
\[
 (Nn+j)!
 =
 j!\,N^{Nn}
 \prod_{r=1}^{N}\left(\frac{j+r}{N}\right)_n.
\]
Exactly one factor in the product is \((1)_n=n!\).  Removing it leaves the
parameter set \(\mathbf b_{N,j}\) and gives the hypergeometric series.
\end{proof}

\begin{corollary}[Factorial-lane genealogy]
\label{cor:block15-factorial-genealogy}
The step-factorial split supplies an exact multiplicative partition of the
factor set, while the residue projectors supply an exact additive partition of
the coefficient series.  The residue projectors instantiate
\cref{thm:block15-universal-lane-transport} directly.  The step-factorial
product is realized by the supplied finite multilinear pointwise
multiplication map on the module of lane sequences, as in
\cref{thm:block15-multilinear-lane-expansion}.  Their common address
refinement is divisibility of the modulus.  Hence the additive transport law
and the multiplicative companion share exact address-refinement
combinatorics, without identifying a product factorization with an additive
idempotent splitting.
\end{corollary}

\subsection{Typed native specializations}
\label{subsec:block15-specializations}

The next statements identify concrete realizations and record the hypotheses
that make each use of the theorem legal.  They do not identify the surrounding
strata with one another.

\subsubsection{Fourier and polyphase lanes}

\begin{proposition}[Typed polyphase realization]
\label{prop:block15-polyphase-specialization}
Let \(K\) be a field, let \(N\in\mathbb Z\) with \(N\ge1\) be invertible in \(K\), and suppose
\(K\) contains a primitive \(N\)-th root of unity.  For
\(F(z)=\sum_{n\ge0}a_nz^n\in K[[z]]\), set
\[
 F_j(w):=\sum_{n\ge0}a_{Nn+j}w^n.
\]
Then
\[
 F(z)=\sum_{j=0}^{N-1}z^jF_j(z^N)
\]
is exact analysis and synthesis.  The Fourier residue projectors and the
polyphase components satisfy the explicit coordinate relation
\[
 \mathcal P_{N,j}F(z)=z^jF_j(z^N).
\]
For every integer \(M\ge1\), refinement from \(N\) to \(NM\) is the
projector-sum law of
\cref{thm:block15-factorial-residue-realization}.
\end{proposition}

\begin{proof}
Every exponent has a unique quotient and remainder on division by \(N\).
The Fourier formula in
\cref{thm:block15-factorial-residue-realization} identifies the same
coefficient sections with character projectors.
\end{proof}

\subsubsection{Cartier and fixed-radix lanes}

\begin{proposition}[Typed Cartier realization]
\label{prop:block15-cartier-specialization}
Let \(K\) be a field and let \(k\ge2\).  Work with \(K\)-valued formal
series and fix the least-significant-digit-first convention.  Define
\[
 \Lambda_d\left(\sum_{n\ge0}a_nz^n\right)
 :=\sum_{n\ge0}a_{kn+d}z^n,
 \qquad
 S_d(G):=z^dG(z^k),
 \quad 0\le d<k.
\]
Then
\[
 \Lambda_dS_e=\delta_{de}1,
 \qquad
 \sum_{d=0}^{k-1}S_d\Lambda_d=1.
\]
For a digit word \(w=d_0d_1\cdots d_{a-1}\), the iterated Cartier lane is
\[
 \Lambda_{d_{a-1}}\cdots\Lambda_{d_1}\Lambda_{d_0}F.
\]
It selects coefficients with index
\[
 k^an+[w]_k,
 \qquad
 [w]_k=d_0+kd_1+\cdots+k^{a-1}d_{a-1}.
\]
Concatenation gives strict radix refinement and the lane path law.
\end{proposition}

\begin{proof}
The first identities follow by coefficient extraction and unique quotient and
remainder.  Iterating in the stated reading convention gives the displayed
address, and concatenation composes the affine index maps.
\end{proof}

\subsubsection{Holonomic lanes}

\begin{proposition}[Typed holonomic realization]
\label{prop:block15-holonomic-specialization}
Let \(K\) be a characteristic-zero field and let the \(K\)-valued sequence
\((u_n)\) satisfy a linear
recurrence with polynomial coefficients.  Beyond a finite singular prefix,
write the recurrence as a finite-dimensional rational nonautonomous state
model
\[
 v_{n+1}=A(n)v_n,
 \qquad A(n)\in\operatorname{Mat}_d(K(n)).
\]
For a fixed integer \(N\ge1\) and residue \(0\le j<N\), the lane state
\(w_m=v_{Nm+j}\) evolves by the exact \(N\)-step rational transfer
\[
 w_{m+1}
 =
 A(Nm+j+N-1)\cdots A(Nm+j)w_m
\]
whenever the displayed factors are defined.  Equivalently, residue
multisection of a D-finite generating series gives D-finite lanes.  The finite
singular prefix and any exceptional denominators remain explicit data.
\end{proposition}

\begin{proof}
Iterate the state equation \(N\) times.  Closure of D-finite series under
scaling the variable, finite sums, and root-of-unity multisection gives the
equivalent generating-series statement after a field extension containing the
needed roots.  No autonomous finite-state claim is used.
\end{proof}

\subsubsection{Local Newton--Mahler lanes}

\begin{proposition}[Typed local clopen realization]
\label{prop:block15-local-specialization}
Fix a prime \(p\) and an integer \(a\ge0\).  For
\(f\in C(\mathbb Z_p,\mathbb Z_p)\), with \(q\in\mathbb Z_p\) the branch
coordinate, define
\[
 \Analysis_{p^a}(f)
 :=
 \bigl(q\mapsto f(p^aq+s)\bigr)_{0\le s<p^a}.
\]
Piecewise synthesis on the disjoint clopen cosets
\(s+p^a\mathbb Z_p\) is its inverse.  Refining from \(p^a\) to
\(p^{a+b}\), for every integer \(b\ge0\), gives the exact child-address and
mixed-radix composition law.
Let
\[
 \mathcal M:=\mathcal M_p^{-1}:C(\mathbb Z_p,\mathbb Z_p)
 \xrightarrow{\sim}c_0(\mathbb Z_p)
\]
be Mahler analysis from \cref{thm:block08-mahler-coordinate}, and write
\(R_{p^a,s}f(q):=f(p^aq+s)\).  Conjugation gives the exact coefficient
operator
\[
 \mathcal T_{p^a,s}
 :=\mathcal M R_{p^a,s}\mathcal M^{-1}:
 c_0(\mathbb Z_p)\longrightarrow c_0(\mathbb Z_p),
\]
whose affine kernel is computed in \cref{thm:block08-affine-transfer}.  Its
mixed-radix composition is \cref{thm:block08-kernel-composition}.

The assertion of complete clopen analysis is restricted to these
prime-power packets.  More general affine coefficient kernels may still
compose according to their declared mixed-radix law, but an arbitrary integer
modulus does not automatically partition \(\mathbb Z_p\) into that many
disjoint clopen branches.
\end{proposition}

\begin{proof}
The \(p^a\) cosets are disjoint, clopen, and cover \(\mathbb Z_p\).  Restricting
after piecewise synthesis and synthesizing after restriction are identities.
Writing a child address as \(s+p^at\) proves refinement.  The cited Mahler
isomorphism makes the displayed conjugation exact, and the two cited kernel
theorems give its coefficient formula and composition law.
\end{proof}

\subsubsection{Nonlinear tree lanes}

\begin{proposition}[Typed Fuss--Catalan realization]
\label{prop:block15-tree-specialization}
Let \(K\) be a characteristic-zero field, let \(q\in\mathbb Z\) with
\(q\ge2\), and let
\(\alpha\in K^\times\).  The uniquely marked formal solution of
\[
 U=z+\alpha U^q,
 \qquad U(0)=0,
\]
has support
\[
 [z^{1+(q-1)n}]U
 =\frac{1}{(q-1)n+1}\binom{qn}{n}\alpha^n
 \ne0
 \qquad(n\ge0),
\]
and every other coefficient is zero.  Thus its support is exactly the set of
degrees congruent to \(1\pmod{q-1}\).
For an integer packet width \(N\ge1\), set
\(g=\gcd(N,q-1)\).  Exactly \(N/g\) residue addresses are active, namely
the congruence class \(1\pmod g\).  Packetizing the specified multilinear
map \((U_1,\ldots,U_q)\mapsto U_1\cdots U_q\) gives the exact finite
carry-convolution system supplied by
\cref{thm:block15-multilinear-lane-expansion}.  On a cyclic orbit of order
\(q-1\), the power map used for packet synthesis has
\[
 \frac{q-1}{\gcd(N,q-1)}
\]
distinct images.  If a fixed embedding \(K\hookrightarrow\mathbb C\) sends
\(\alpha\) to a positive real number and the marked branch is interpreted in
the smooth implicit analytic schema, its critical points and values satisfy
\[
 q\alpha\tau^{q-1}=1,
 \qquad
 \rho=\tau-\alpha\tau^q=\frac{q-1}{q}\tau.
\]
The resulting orbit \(\rho\mu_{q-1}\) is precisely the dominant singularity
orbit, and packetization leaves \((q-1)/g\) distinct images.
\end{proposition}

\begin{proof}
Lagrange inversion gives the displayed coefficient formula and zero at every
other degree.  Characteristic zero and \(\alpha\ne0\) make every displayed
coefficient nonzero.  Their residues modulo \(N\) form the coset of the subgroup
generated by \(q-1\), whose size is \(N/g\).  Multilinear lane expansion of
the product gives the packet equation, including the power of the packet
variable produced by the quotient in residue addition.  On the cyclic group
of order \(q-1\), the kernel of multiplication by \(N\) has order \(g\), so
the image has order \((q-1)/g\).  The critical equations give the displayed
values, while \cref{thm:block11-dominant-orbit} identifies their full orbit as
the dominant one under the stated analytic hypotheses.
\end{proof}

\subsubsection{Algebraic branch lanes}

\begin{proposition}[Typed cyclic-radical realization]
\label{prop:block15-algebraic-specialization}
Let \(K\) be a characteristic-zero field containing \(\mu_d\) and a
primitive \(N\)-th root of unity \(\omega\), let
\(d,m,N\in\mathbb Z_{\ge1}\), and choose
the formal branch
\[
 F(z)=(1-\alpha z^m)^{-1/d},
 \qquad \alpha\in K^\times.
\]
With \(g=\gcd(N,m)\), its nonzero coefficients occupy exactly the subgroup of
\(\mathbb Z/N\mathbb Z\) generated by \(m\), which has \(N/g\) addresses.
The rotated radicands
\[
 1-\alpha\omega^{jm}z^m,
 \qquad 0\le j<N,
\]
also have exactly \(N/g\) distinct values.

If, in addition, their distinct Kummer classes generate a subgroup isomorphic
to \((\mathbb Z/d\mathbb Z)^{N/g}\) in the Kummer quotient
\[
 K(z)^\times/K(z)^{\times d},
\]
then the compositum of the corresponding radical branch fields has degree
\[
 d^{N/g}
\]
over \(K(z)\).  Without that independence hypothesis, the lane orbit count
remains exact but the degree formula is only an upper bound.
\end{proposition}

\begin{proof}
The binomial series is supported at exponents \(mn\).  The cyclic subgroup
calculation gives both address and rotation counts.  The degree statement is
the standard Kummer independence calculation for the explicitly listed
classes.  A relation among the classes produces a field collision not detected
by address count alone.
\end{proof}

\subsubsection{Finite Galois lanes}

\begin{proposition}[Typed semisimple Galois realization]
\label{prop:block15-galois-specialization}
Let \(G\) be finite and let \(K\) be a characteristic-zero splitting field
for \(G\).  The primitive central idempotents
\[
 e_\chi
 =
 \frac{\chi(1)}{|G|}
 \sum_{g\in G}\chi(g^{-1})g
 \in K[G]
\]
for \(\chi\in\operatorname{Irr}_K(G)\) form a complete
pairwise-annihilating lane system on every finite-dimensional
semisimple \(K[G]\)-module \(V\):
\[
 V=\bigoplus_{\chi\in\operatorname{Irr}_K(G)} e_\chi V.
\]
Every \(G\)-equivariant map is block diagonal in these isotypic lanes, while
an arbitrary \(K\)-linear map has the full lane matrix of
\cref{thm:block15-universal-lane-transport}.

For a subgroup \(H\le G\), the Reynolds operator
\[
 e_H=\frac1{|H|}\sum_{h\in H}h
\]
is an idempotent projector onto \(V^H\); equivalently,
\((e_H,1-e_H)\) is a two-lane system on the underlying \(K\)-vector space.
It is not in general a Boolean
selector formed from the isotypic projectors above.  Projectors belonging to unrelated
subgroups need not annihilate one another, so they do not automatically form
one lane system.  Transport through induction, restriction, or subgroup
refinement requires the corresponding Mackey or double-coset data.
\end{proposition}

\begin{proof}
Maschke semisimplicity and the character idempotent relations give the first
decomposition.  An equivariant map commutes with the center of \(K[G]\), hence
with every \(e_\chi\).  Averaging proves that \(e_H\) is idempotent with image
\(V^H\).  Two subgroup averages generally fail the pairwise-annihilation
axiom, which proves the stated boundary.
\end{proof}

\subsection{Descent of lane systems through a finite atlas}
\label{subsec:block15-lane-descent}

The local theorem becomes global only when the addresses and splittings glue
without address holonomy.  The following statement makes those additional
hypotheses explicit.

\begin{theorem}[Finite atlas descent for lane transport]
\label{thm:block15-lane-descent}
Let \(\mathfrak R\) be a finite profiled atlas whose participating chart
categories are additive and idempotent-complete, with additive transition
functors \(F_{UV}:\mathfrak R(U)\to\mathfrak R(V)\) on its declared
overlaps.  Let
\[
 X=(X_U,\theta^X_{UV}),
 \qquad
 Y=(Y_U,\theta^Y_{UV}),
 \qquad
 Z=(Z_U,\theta^Z_{UV})
\]
be three compatible descent objects.  Thus the transition isomorphisms satisfy
the identity and cocycle equations.  Let
\[
 f=(f_U):X\longrightarrow Y,
 \qquad
 g=(g_U):Y\longrightarrow Z
\]
be compatible descended morphisms.

Fix one finite resolution shape \(\mathsf R\).  Suppose every local object
\(A_U\), for \(A\in\{X,Y,Z\}\), carries a finite profiled lane system over
\(\mathsf R\).  For every overlap, every \(r\in\mathsf R\), and every such
\(A\), suppose there is an address bijection
\[
 \lambda_{UV}^{A,r}:
 I_{A_U,r}\xrightarrow{\sim}I_{A_V,r}
\]
such that
\[
 \lambda_{UU}^{A,r}=1,
 \qquad
 \lambda_{UT}^{A,r}
 =\lambda_{VT}^{A,r}\lambda_{UV}^{A,r},
\]
on every triple overlap \((U,V,T)\), and, for \(s\succeq r\),
\[
 \rho_{sr}^{A_V}\lambda_{UV}^{A,s}
 =
 \lambda_{UV}^{A,r}\rho_{sr}^{A_U}.
\]
Require trivial address holonomy: for every \(A\) and \(r\), there is a
global finite address set \(I_{A,r}\), a global refinement map
\(\rho_{sr}^{A}:I_{A,s}\twoheadrightarrow I_{A,r}\), and bijections
\[
 \mu_U^{A,r}:I_{A_U,r}\xrightarrow{\sim}I_{A,r}
\]
such that
\[
 \mu_V^{A,r}\lambda_{UV}^{A,r}=\mu_U^{A,r},
 \qquad
 \mu_U^{A,r}\rho_{sr}^{A_U}
 =\rho_{sr}^{A}\mu_U^{A,s}.
\]
Assume also that the transition isomorphisms intertwine the local projectors:
\[
 \theta^A_{UV}
 F_{UV}(p_{A_U,i}^r)
 =
 p_{A_V,\lambda_{UV}^{A,r}(i)}^r
 \theta^A_{UV},
\]
and that the verified transition profiles preserve the chosen splittings,
finite sums, composition, and refinement equations, with coherent splitting
comparisons on triple overlaps.

Then:
\begin{enumerate}[label=\textup{(D\arabic*)}]
\item after relabeling by the \(\mu_U^{A,r}\), each local lane and its chosen
splitting define a lane summand in \(\Desc(\mathfrak R)\), and the global
addresses carry the declared projector-sum refinement;
\item the local matrices \(([f_U]_{a i})\), \(([g_U]_{c a})\), and
\(([g_Uf_U]_{c i})\) glue after the prescribed address relabelings and
splitting comparisons;
\item the descended matrices satisfy
\[
 [gf]_{c i}=\sum_a[g]_{c a}[f]_{a i};
\]
\item every compatible family of support selectors and named residuals
descends, and every compatible refinement square remains a projector-sum
refinement square in the descent category;
\item every fold whose source objects, target objects, and structure maps are
compatible descent data defines a fold in \(\Desc(\mathfrak R)\).  Compatible
selectors and left inverses descend, so a split-losslessness certificate
descends directly from its left-inverse identity.  Moreover, suppose the
relevant local kernels \(K_U=\ker(\Phi_U)\) exist, every transition functor
preserves them, and there are coherent comparison isomorphisms
\[
 \eta_{UV}:F_{UV}(K_U)\xrightarrow{\sim}K_V
\]
compatible with the kernel inclusions and satisfying the triple-overlap
cocycle.  Then the global kernel exists and is computed chartwise:
\[
 \bigl(\ker_{\Desc(\mathfrak R)}\Phi\bigr)_U
 \cong
 \ker_{\mathfrak R(U)}(\Phi_U).
\]
\end{enumerate}
\end{theorem}

\begin{proof}
For a global address \(i\in I_{A,r}\), use the local projector indexed by
\((\mu_U^{A,r})^{-1}(i)\).  The trivialization identity makes these projectors
correspond under the transition isomorphisms on every overlap; its refinement
identity makes their local
projector-sum laws agree.  The coherent splitting comparisons make the
corresponding inclusions and projections descent morphisms.  Finite
biproducts in the descent category are computed chartwise, so local analysis
and synthesis glue and remain mutually inverse.

Compatibility of \(f\) and \(g\), followed by the splitting comparisons,
shows that their local coordinate blocks agree on overlaps after applying the
specified address relabelings.  Hence the matrices descend.  Matrix
multiplication is a finite sum of compositions, both computed chartwise, so
the local path law glues to the displayed global law.  The same argument
applies to projector sums, selectors, and residual differences.

Finally, a compatible fold is a morphism in the relevant descent category.
Compatible selectors and left inverses glue, and their identities can be
checked chartwise.  Under the additional kernel hypotheses, the objects
\((K_U,\eta_{UV})\) form a descent object; compatibility with the local kernel
inclusions gives a global morphism into the fold source.  Its universal
property is checked in each chart, which proves the displayed chartwise
kernel identity.  The split case needs no kernel-preservation hypothesis:
the compatible local left inverses glue directly.
\end{proof}

\begin{corollary}[Certified atlas lane transport]
\label{cor:block15-certified-atlas-lane-transport}
Let \(\mathfrak A_{\mathrm{cert}}\) be the certified finite atlas of this
manuscript.  On every subatlas whose passports include additive split-lane
systems, compatible refinement maps, coherent address relabelings with
trivial address holonomy, and the required preservation profile, the
Universal Lane Transport Theorem holds in the descent category.  Kernel
claims additionally require the kernel-preservation and coherence data in
\cref{thm:block15-lane-descent}.  In particular, any of the typed
specializations above may be transported across a certified reversible
overlap after its native hypotheses, address trivializations, and splitting
cocycles have been checked.
\end{corollary}

\begin{proof}
Apply \cref{thm:block15-lane-descent} to the verified subatlas.  The conclusion
is deliberately restricted to the overlaps whose passports contain all of
the listed data.
\end{proof}

\begin{boundary}[Descent is conditional on address coherence]
\label{bnd:block15-descent-boundary}
An equivalence of underlying semantic objects does not by itself identify two
lane systems.  Without the bijections \(\lambda_{UV}^{A,r}\), their cocycle
law, compatible global trivializations \(\mu_U^{A,r}\), refinement
compatibility, and coherent splitting comparisons, the local matrices need
not glue as ordinary globally addressed lanes.  Nontrivial address holonomy
can permute lanes around a loop; then only invariant orbit sums, or lanes in a
separately defined local-system-valued category, may descend.  The theorem
above verifies ordinary descent under trivial holonomy.  It does not
manufacture missing lane passports for an arbitrary heterogeneous atlas.
\end{boundary}

\subsection{Sharp boundaries and a finite counterexample}
\label{subsec:block15-boundaries}

\begin{proposition}[Endpoint selection does not implement refinement]
\label{prop:block15-endpoint-counterexample}
Let \(K\) be a field and let \(X=K^2\).  At the coarse resolution use the
single projector \(p=1_X\).  At the fine resolution use
\[
 p_1=
 \begin{pmatrix}1&0\\0&0\end{pmatrix},
 \qquad
 p_2=
 \begin{pmatrix}0&0\\0&1\end{pmatrix},
 \qquad p=p_1+p_2.
\]
For
\[
 f=p_2,
\]
the unique coarse block is \(pfp=f\ne0\), whereas selecting the single fine
endpoint \((p_1,p_1)\) gives \(p_1fp_1=0\).  Therefore a coarse block is the
sum over all compatible descendant source-target pairs.  It cannot in general
be represented by choosing one fine endpoint.
\end{proposition}

\begin{proof}
The displayed matrix multiplication gives the two values immediately.  The
full refinement identity recovers
\(f=\sum_{a,i\in\{1,2\}}p_afp_i\), whose only nonzero term here is
\(p_2fp_2\).
\end{proof}

\begin{boundary}[Finite profiled meaning of ``universal'']
\label{bnd:block15-universal-boundary}
The theorem proves a common calculus for finite split lanes in additive
categories and for the explicitly stated profile extensions.  It does not
assert any of the following:
\begin{enumerate}
\item that every mathematical object has a canonical or minimal lane system;
\item that countably infinite lane sums converge without a topology and a
convergence theorem;
\item that arbitrary nonlinear, probabilistic, operator-algebraic, or
homotopical constructions preserve additive lane identities;
\item that address count equals vector-space dimension, module rank,
information, entropy, algebraic degree, or computational complexity;
\item that a support pattern, heatmap, norm profile, valuation profile, or
orbit count is a complete invariant;
\item that native invariants in different strata are equal merely because
their formulas contain the same gcd or residue operation;
\item that the theorem supplies historical novelty evidence or external
validation.
\end{enumerate}
Each stronger conclusion requires a native theorem and a preservation
certificate.
\end{boundary}

\begin{remark}[Classical core and present synthesis]
\label{rem:block15-classical-core}
Finite biproduct coordinates, Peirce decompositions, block-matrix
multiplication, residue filters, and polyphase identities have classical
antecedents.  The theorem packages that core with explicit address refinement,
profile-controlled realization, named residuals, exact fold kernels, and
typed descent.  This manuscript treats the resulting assembly as an internal
mathematical contribution.  It does not infer a priority claim without a
separate literature study.
\end{remark}

\subsection{Exact finite validation and interpretation}
\label{subsec:block15-validation}

The accompanying verifier is an exact finite regression suite for selected
implementation clauses.  It checks integer and rational block reconstruction,
matrix composition, refinement aggregation, support and residual identities,
fold collision kernels, factorial-lane partitions, residue-projector
refinement, and explicit negative controls.  The proofs above do not depend on
finite sampling.  Conversely, the symbolic proofs do not certify that a
particular program implements the formulas correctly.  The two records serve
different purposes.

\IfFileExists{generated/block15_validation_table.tex}{%
\input{generated/block15_validation_table.tex}%
}{%
\begin{table}[htbp]
\centering
\caption{Exact finite regression for selected Universal Lane Transport
implementation clauses.  The generated validation table is inserted when the
verifier has been run.}
\label{tab:block15-lane-transport-validation-placeholder}
\begin{tabular}{@{}ll@{}}
\toprule
Record & Status\\
\midrule
Machine-readable validation & Run the Block 15 verifier\\
\bottomrule
\end{tabular}
\end{table}%
}

\begin{resultbox}
\textbf{Universal Lane Transport result.}
Finite split lane systems admit exact analysis and synthesis, unique transfer
matrices, path composition, strict projector-sum refinement, Boolean support
selectors, named residuals, and exact fold-loss kernels.  Complete changes of
resolution are coordinate equivalences.  In abelian targets, fold loss sits in
a canonical short exact sequence.  In module targets, specified multilinear
and finite polynomial operations expand lane by lane.  Residue/Fourier,
polyphase, Cartier, holonomic, local Newton--Mahler, Fuss--Catalan, cyclic
radical, and semisimple Galois systems realize these clauses under their
printed hypotheses.  The step-factorial product is an exact multiplicative
companion transported through a specified multilinear multiplication map,
not an additive idempotent splitting.  Compatible address trivializations,
refinement, and splitting data make the calculus descend through a finite
certified atlas when address holonomy is trivial.  No claim is made beyond
that finite profiled domain.
\end{resultbox}
